% SPDX-License-Identifier: CC-BY-4.0
% Copyright (c) 2026 Martin Schr\"oder <info@swedishembedded.com>

\begin{abstract}
A language model's hidden state could be stored as a few Fourier coefficients
of a longer virtual signal and processed by per-frequency operators, cutting
parameters and arithmetic. We first show that a basis change alone cannot do
this: a decoder whose maps on the residual stream are dense absorbs any
orthogonal basis into its weights, and truncating its spectrum yields an
ordinary narrow decoder. Consistent with this, the lowest Fourier modes of
Qwen3-0.6B's residual stream keep about as much of it as random directions
do, while its principal components keep nearly all. We then define
SpectralLM, a decoder whose residual stream is $m$ Fourier coefficients and
whose feed-forward networks are spectral mixers: per-mode filters around a
gated product taken in the signal domain, $3Cm$ parameters per block. We
implement it with gradient-checked GPU kernels, a fast-transform path held to
a dense reference, and incremental decoding, and train it from scratch on
TinyStories against Qwen3-architecture controls of the same width. At equal
parameters SpectralLM reaches \BpbS{} bits per byte on held-out stories
against \BpbQpar{} for the dense control (difference \DiffSQpar{}, 95\%
interval \DiffSQparCI{}, three seeds each), so it is non-inferior; at about
equal arithmetic a dense feed-forward network is better by \DiffSQstd{} bits
per byte, and SpectralLM trains at half its speed with several times its
memory. The pretrained Qwen3-0.6B scores \BpbQwen{} zero-shot and
\BpbQwenFt{} after 4.2M tokens of continued training. At this scale the
spectral mixer buys parameter efficiency, not compute efficiency.
\StatusAbstract
\end{abstract}
